\documentclass[10pt,a4paper]{amsart}
\usepackage[margin=19mm]{geometry}
\usepackage{amsmath,amssymb,mathtools}
\usepackage[scr=rsfs,cal=euler]{mathalfa}
\usepackage{xfrac}
\usepackage[unicode,hidelinks,bookmarks=false]{hyperref}
\newcommand{\ie}{i.e.\ }
\newcommand{\cf}{cf.\ }
\newcommand{\resp}{respectively\ }
\newcommand{\Subsection}{\subsection}
\providecommand{\nobreakdash}{\nobreak\hbox{-}}
\setlength{\emergencystretch}{2em}
\multlinegap0pt
\usepackage{amssymb}
\usepackage{enumerate}
\usepackage[all]{xy}
\usepackage{tikz}
\usepackage[shortlabels]{enumitem}
\newtheorem{lem}{Lemma}
\title{On the source algebra equivalence class of blocks with cyclic defect groups, III}
\author{Gerhard Hiss, Caroline Lassueur}
\date{Source: arXiv:2601.01582v1 (CC BY 4.0)}
\begin{document}
\maketitle

\noindent\emph{Source and licence.} On the source algebra equivalence class of blocks with cyclic defect groups, III. Version arXiv:2601.01582v1 (CC BY 4.0), distributed by arXiv under the Creative Commons Attribution 4.0 International licence, http://creativecommons.org/licenses/by/4.0/, as recorded in the arXiv licence metadata for this exact version and shown on the abstract page https://arxiv.org/abs/2601.01582v1. Full licence notice: \texttt{licenses/arxiv-2601.01582-CC-BY-4.0.txt}.\par\smallskip

\noindent\textit{Editorial scope.} Counted unit: the article's lem YEqualOpZ (source0.tex lines 597--615). The article's complete original proof of it is delivered with it (source0.tex lines 616--632, one \texttt{proof} environment of 17 lines). No mathematics is written, repaired or re-derived: the statement and the proof are the article's own lines, in the article's own notation, and no retained line is substituted or re-flowed.  No further source-local context is needed: every cross-reference of every retained line resolves inside the two delivered spans.  Nothing was omitted from any retained span, so the package carries no omission record.  No retained line contains a citation, so the wrapper carries no bibliography and no bibliography entry.  No retained line contains a figure environment, an image-inclusion command or drawing code, so no image is delivered and the package declares no asset dependency.  Editorial formatting only, in addition: an AMS \texttt{amsart} preamble that restates the article's own declarations and macros used by the retained lines, namely the environments \texttt{lem} on separate counters, together with the standard packages those lines need.  The retained environments print in this document as Lemma 1 in order of appearance, whereas the article prints the same items with the numbering of its own full document.  The complete native source of the article as distributed in the arXiv e-print for this exact version is bundled unchanged as \texttt{sources/192089/source0.tex}.
\begin{lem}
\label{YEqualOpZ}
{\rm (a)} We have $Y \leq Z( \tilde{G} )$.

{\rm (b)} The kernel of the natural epimorphism 
$\langle t \rangle \rightarrow \bar{D}$ equals $\langle t \rangle \cap Y$, so
that $|\bar{D}| = |t|/p^{c'}$.

{\rm (c)} We have $O_p( Z ) = Z \cap D$. Moreover, 
$D = \langle t, Y \rangle = \langle t, O_p(Z) \rangle$. In particular, 
$r(D) \leq 2$.

{\rm (d)} If $\langle t \rangle \cap O_p( Z ) = \{ 1 \}$, then $Y = O_p( Z )$.

{\rm (e)} If $|D| = |t|p^{a}$ then 
$\langle t \rangle \cap O_p( Z ) = \{ 1 \}$ and $p^a \mid n$.

{\rm (f)} If~$D$ is cyclic, then $D = \langle t \rangle$.
\end{lem}

\begin{proof}
(a) and (b) are trivial.

(c) Since $D$ is a defect group, we have $O_p( Z ) \leq D$, implying the first 
assertion. The other assertions follow
from $\bar{D} = \langle \bar{t} \rangle$ and $Y \leq O_p(Z)$.

(d) By~(c) and the assumption, $D = \langle t \rangle \times O_p( Z )$. Since 
$Y \leq O_p( Z )$, the quotient $D/Y$ can only be cyclic if $Y = O_p( Z )$.

(e) We have $|t|p^a = |D| \leq |t||Y| \leq |t||O_p(Z)| \leq |t|p^a$, implying
the claims.

(f) If~$D$ is cyclic and~$\bar{D}$ is non-trivial, the image of a proper subgroup 
of~$D$ is a proper subgroup of~$\bar{D}$. As the image of $\langle t \rangle$
equals~$\bar{D}$, we obtain our claim.
\end{proof}

\end{document}
